\begin{table}[htbp]\centering\small
\caption{The reference and the data}\label{tab:fit}
\begin{tabularx}{\linewidth}{@{}>{\raggedright\arraybackslash}Xrr@{}}\toprule
 & Model & Data \\\midrule
Private employment (million) & 117.9 & 118.3 \\
Knowledge share of private employment (percent) & 64.8 & 63.2 \\
Private unemployment rate (percent) & 2.1 & 1.9 \\
Inactive residents (million) & 52.2 & 52.7 \\
Retired residents (million) & 47.7 & 47.5 \\
Other employed residents (million) & 44.6 & 44.2 \\
Unemployed knowledge workers taking manual jobs (percent per quarter) & 10.1 & 10.4 \\
Unemployed manual workers taking knowledge jobs (percent per quarter) & 10.2 & 9.3 \\
Employed knowledge workers moving to manual jobs (percent per quarter) & 0.21 & 0.68 \\
Employed manual workers moving to knowledge jobs (percent per quarter) & 0.86 & 0.96 \\
Residents without liquid assets (percent) & 23.1 & 14.1 \\
Liquid assets (quarters of GDP) & 1.32 & 1.20 \\
\bottomrule\end{tabularx}
\par\smallskip\begin{minipage}{\linewidth}\footnotesize\emph{Notes.} Data moments are those of the household calibration described in the text and in Appendix~\ref{sec:data}. Liquid assets are divided by quarterly value added of the model in both columns. The four flows are not the same object in the two columns. In the model a flow is the share of the residents in the origin status at the start of a quarter who hold a job in the other occupation when income is paid in that quarter. In the data it is a monthly rate $r$ converted as $1-(1-r)^3$, the cumulative probability over a quarter under a constant monthly hazard. With other destinations and with returns, this can differ from the probability of being in the other occupation three months later, which has not been estimated directly from the linked records.\end{minipage}
\end{table}
